%  =============================================================================
%  -> Abschnitt "kapitel/03_umsetzung.tex"
%
%  Gliederung an der KI-Fallstudie ausgerichtet:
%  Architektur -> Strategie -> Implementierung -> Evaluation -> Wuerdigung.
%  =============================================================================

\chapter{Fallstudie: OpenClaw als Trading Agent}
\label{cha:umsetzung}

Nachdem das vorangegangene Kapitel mit dem Periodensystem der Künstlichen Intelligenz
einen Ordnungsrahmen bereitgestellt hat, überträgt das folgende Kapitel diesen Rahmen
auf den konkreten Gegenstand der Arbeit. Der untersuchte Agent verfolgt eine auf
gleitenden Durchschnitten beruhende Handelsstrategie und löst die daraus abgeleiteten
Entscheidungen eigenständig als Ordervorgänge über eine Bankschnittstelle aus. Damit
schließt er den eingangs beschriebenen Kreis aus Wahrnehmung, Schlussfolgerung und
Handlung vollständig und lässt sich als Verbindung mehrerer KI-Elemente im Sinne des
Periodensystems auffassen.

Auf der Ebene der Wahrnehmung stützt sich der Agent auf die kontinuierliche Erfassung
und Aufbereitung von Marktdaten, deren Qualität über eine unabhängige Zweitquelle
gegengeprüft wird, sowie auf die Erkennung von Trendmustern in der Kurszeitreihe. Auf
der Ebene der Schlussfolgerung tritt ein Gremium aus drei großen Sprachmodellen als
bewertende Instanz hinzu, das die regelbasiert vorbereitete Entscheidung jeweils
unabhängig prüft, mit einer nachvollziehbaren Begründung versieht und dessen Voten nach
einer strikten Mehrheitsregel aggregiert werden; deterministische Vorprüfungen
entscheiden harte Ausschlusskriterien bereits vor dem ersten Modellaufruf. Die
Modellinferenz erfolgt über eine einheitliche Schnittstelle, die den Wechsel zwischen
cloud-gehosteter Ausführung und dem lokalen Betrieb auf eigener Hardware als reine
Konfigurationsentscheidung belässt. Auf der Ebene der Reaktion schließlich ordnet
sich das Vorhaben unmittelbar dem Element der prozesssteuernden Kontrolle zu, das der
zugrunde gelegte Rahmen ausdrücklich mit dem Börsenhandel als Anwendung
verbindet.\zit[S. 18]{waldhoerKI2025} Ergänzend sorgt eine Kommunikationskomponente für
die Rückbindung an den Menschen: Sie übermittelt die tägliche Freigabe der Banksitzung,
berichtet jeden Entscheidungslauf und stellt einen jederzeit wirksamen Not-Aus bereit;
ein öffentliches, redigiertes Dashboard macht die Entscheidungs- und Wertentwicklung
zusätzlich transparent.

Kennzeichnend für die gewählte Gestaltung ist weniger die Breite der eingesetzten
Fähigkeiten als vielmehr deren bewusste Begrenzung. Die Sprachmodelle werden nicht zur
Prognose künftiger Kurse und nicht zur erfahrungsbasierten Selbstveränderung eingesetzt;
die entsprechenden Elemente des Periodensystems bleiben bewusst unbesetzt, um die
Nachvollziehbarkeit und Reproduzierbarkeit des Systemverhaltens zu wahren. Das Gremium
kann Handelsvorschläge ausschließlich verhindern, verkleinern oder zeitlich einordnen,
niemals erzeugen oder in ihrer Richtung umkehren. Flankiert wird diese begrenzte
Diskretion durch einen mehrstufig abgesicherten Orderpfad: Ein zweiphasiges
Ausführungsjournal verhindert die Doppelausführung desselben Entscheids, ein
Orderbuchabgleich klärt unklare Ausführungszustände, die Scharfschaltung auf reale
Orders erfolgt ausschließlich kontospezifisch, und jeder unklare Zustand endet in einem
alarmierten Handelsstopp statt in einem automatischen Wiederholungsversuch. Die
eigentliche gestalterische Aussage der Arbeit liegt somit in der sicherheitsorientierten
Einschränkung des KI-Einsatzes, deren Ausprägung im weiteren Verlauf des Kapitels
dargelegt wird.

Der Aufbau des Kapitels folgt dem Weg von der Struktur zur Bewertung. Zunächst wird die
Systemarchitektur des Agenten mit ihren Komponenten und Datenflüssen beschrieben. Darauf
aufbauend wird die Handelsstrategie formal gefasst und ihre Umsetzung in Quelltext
dargestellt. Die anschließende Evaluation ordnet die Ergebnisse anhand geeigneter
Kennzahlen ein und benennt die methodischen Grenzen der Untersuchung. Eine kritische
Würdigung der mit dem autonomen Handel verbundenen Risiken und der vorgesehenen
Absicherungsmechanismen beschließt das Kapitel.

\section{Systemarchitektur des Agenten}
\label{sec:umsetzung-architektur}

Die Architektur des Agenten folgt einem durchgängigen Gestaltungsprinzip: Module
kommunizieren nicht durch direkte Aufrufe, sondern ausschließlich über gemeinsamen,
persistierten Zustand in Form von SQLite-Datenbanken und synchronisierten Dateien.
Diese lose Kopplung erlaubt es, jede Komponente isoliert zu testen, auszutauschen und
im Fehlerfall neu zu starten, ohne dass sich Störungen über Prozessgrenzen hinweg
fortpflanzen. \Cref{fig:zielarchitektur} zeigt die Zielarchitektur mit den bereits
umgesetzten und den geplanten Bestandteilen.

% Grafik: abbildungen/zielarchitektur.svg — Einbindung als Vektorgrafik via
% \includesvg (svg-Paket, \svgpath{{abbildungen/}}, analog fom-logo).
% Das SVG ist attributbasiert gestylt (kein <style>-Block), damit die
% Inkscape-Konvertierung von Overleaf verlustfrei arbeitet.
\begin{figure}[htbp]
  \centering
  \includesvg[width=\textwidth]{zielarchitektur}
  \caption[Zielarchitektur des autonomen Handelsagenten]{Zielarchitektur des
    autonomen Handelsagenten. Durchgezogene Elemente sind umgesetzt und in Betrieb,
    gestrichelte Elemente geplant.\quelleeigenki}
  \label{fig:zielarchitektur}
\end{figure}

Das produktive System verteilt sich auf zwei virtuelle Maschinen der Google Cloud
Platform. Die \emph{collector-vm} sammelt rund um die Uhr Marktdaten des S\&P~500:
Rohticks werden über einen WebSocket-Datenstrom sowie ein ergänzendes Polling für
Futures-Randzeiten erfasst, durch einen Sanitätsfilter geprüft und zu Minuten- und
Tagesbars verdichtet. Der jeweils jüngste Tagesschluss wird gegen die offizielle
Zweitquelle der Federal Reserve Bank of St.\ Louis (FRED) abgeglichen; Divergenzen
oberhalb einer Schwelle lösen einen Alarm aus. Die \emph{agent-vm} erhält die
Datenbasis als minütlich synchronisierte, atomar ersetzte Replik und beherbergt den
eigentlichen Entscheidungs- und Orderpfad: die deterministische Strategieauswertung,
den mehrstufigen Sprachmodell-Audit, die Orchestrierung der Orderausführung sowie die
Anbindung an die Bankschnittstelle. Beide Maschinen sind wechselseitig überwacht;
ein Totmannschalter auf der collector-vm alarmiert den Betreiber über den
Messaging-Kanal, sobald der Lebenszeichen-Stempel der agent-vm veraltet.

Für den Dauerbetrieb ist die Datensammlung gezielt gehärtet. Bricht der
WebSocket-Datenstrom still ab, erkennt ein Wächter die ausbleibenden Ticks und
erzwingt eine Neuverbindung mit ansteigenden Wartezeiten; der Polling-Prozess für
die Futures-Randzeiten startet sich nach Ausnahmen selbst neu, und ein geordneter
Shutdown stellt sicher, dass keine halb geschriebenen Datensätze zurückbleiben.
Ein Plattenplatz-Wächter, eine wöchentliche Integritätsprüfung der Datenbank und
eine kontinuierliche Replikation des Datenbestands in einen Cloud-Speicher ergänzen
die Vorsorge; Alarme beider Maschinen laufen über eine gemeinsame Brücke im
Messaging-Kanal zusammen. Diese Betriebsdetails sind für die Fragestellung der
Arbeit nicht nebensächlich: Ein Agent, der täglich unbeaufsichtigt Geld bewegen
soll, ist nur so vertrauenswürdig wie die Infrastruktur, die seine Datenbasis und
seine Erreichbarkeit garantiert.

Die Sprachmodellinferenz ist bewusst hinter einer einheitlichen Schnittstelle
(Ollama) gekapselt. Im gegenwärtigen Betrieb werden drei cloud-gehostete Modelle
unterschiedlicher Anbieter befragt; die Migration auf lokal betriebene Modellvarianten
eines Mac~mini ist als reiner Konfigurationswechsel angelegt, da Modellbezeichner,
Endpunkt und Zeitbudgets in der zentralen Konfigurationsdatei liegen und kein
Quelltext angepasst werden muss. Damit bleibt das in \cref{sec:zielsetzung}
formulierte Ziel der lokalen, weitgehend kostenfreien Betreibbarkeit strukturell
gesichert, während der laufende Betrieb von der Verfügbarkeit leistungsfähiger
gehosteter Modelle profitiert.

Eine dritte, strikt entkoppelte Zone bildet das öffentliche Dashboard. Die
collector-vm exportiert minütlich redigierte JSON-Feeds des Entscheidungsjournals
-- Signale, Voten, Orders und Depotwertverlauf, jedoch keinerlei Konto- oder
Zugangsdaten -- in einen privaten Cloud-Storage-Bucket; ein zustandsloser
Cloud-Run-Dienst rendert daraus die Weboberfläche. Die handelnden Maschinen selbst
bleiben ohne offenen eingehenden Port. Im Ordnungsrahmen des Periodensystems der
Künstlichen Intelligenz verbindet die Architektur damit Elemente der Wahrnehmung
(Daten- und Mustererfassung), der Schlussfolgerung (Entscheiden und Erklären) und
der Reaktion, deren Zentrum das Element der prozesssteuernden Kontrolle
bildet.\zit[S.~18]{waldhoerKI2025}

\section{SMA-Handelsstrategie}
\label{sec:umsetzung-strategie}

Die Handelsentscheidungen des Agenten entstammen einer deterministischen
Regimestrategie auf Basis gleitender Durchschnitte. Für die Tagesschlusskurse
$P(t)$ des S\&P~500 ist der einfache gleitende Durchschnitt definiert als
\begin{equation}
  \mathrm{SMA}_n(t) \;=\; \frac{1}{n}\sum_{i=0}^{n-1} P(t-i),
  \qquad n \in \{50,\, 200\}.
  \label{eq:sma}
\end{equation}
Das Handelsregime ergibt sich aus dem Verhältnis der schnellen zur langsamen
Durchschnittslinie:
\begin{equation}
  R(t) \;=\;
  \begin{cases}
    \text{LONG},  & \text{falls } \mathrm{SMA}_{50}(t) > \mathrm{SMA}_{200}(t),\\
    \text{SHORT}, & \text{sonst.}
  \end{cases}
  \label{eq:regime}
\end{equation}
Die Entscheidung für den Handelstag $t{+}1$ nutzt ausschließlich Schlusskurse bis
einschließlich Tag $t$, sodass kein Blick in die Zukunft (Look-Ahead) möglich ist.
Das Regime wird über eine zentrale Zuordnungsdatei auf konkrete Handelsprodukte
abgebildet -- ein Faktor-Zertifikat mit zweifachem Hebel in Long- beziehungsweise
Short-Richtung --, wobei ein Katalog mit Eignungsregeln (Mindestkurs, maximaler
Spread, Hebelobergrenze) die Produktauswahl begrenzt.

Der Kapitaleinsatz je Handelstag wird volatilitätsabhängig skaliert:
\begin{equation}
  S(t) \;=\; \operatorname{clamp}\!\Bigl(K \cdot \min\bigl(1,\,
  \sigma^{*}/\sigma_{20}(t)\bigr),\; S_{\min},\; K\Bigr),
  \label{eq:sizing}
\end{equation}
mit dem Basiskapital $K = 2.000\,€$, der realisierten
20-Tage-Volatilität $\sigma_{20}$ der Index-Tagesrenditen, der Zielvolatilität
$\sigma^{*} = 1{,}0\,\%$ pro Tag und der Mindestgröße $S_{\min} = 1.000\,€$. Die
Obergrenze $K$ ist als harter, weder über die Konfiguration noch durch die
Sprachmodelle veränderbarer Deckel implementiert.

Die Ausführungsregeln sind ebenfalls deterministisch festgelegt. Gehandelt wird im
Fenster der US-Kassasitzung mit drei diskreten Einstiegsoptionen, unter denen das
Prüfgremium wählen darf: dem Einstieg kurz nach der Eröffnungsauktion, einem
bestätigungsbasierten Einstieg nach den ersten fünfzehn Handelsminuten oder einem
bewusst verzögerten Einstieg, der die Eröffnungsvolatilität meidet. Der Ausstieg
folgt dem zuerst erreichten Kriterium aus Verluststopp (drei Prozent), Gewinnziel
(zwei Prozent) und Fensterschluss. Es gilt ein hartes Tagesbudget von höchstens
einem Kauf und einem Verkauf, und sämtliche Orders werden ausschließlich limitiert
aufgegeben, da die Datenfrische im verteilten Aufbau eine Slippage-Kontrolle über
den Limitpreis erfordert.

Die ökonomische Hürde dieser Taktung lässt sich vorab beziffern: Ein Round-Trip
kostet im Direkthandel 7{,}80\,€ an Fixgebühren, hinzu kommt der Geld-Brief-Spread
des Zertifikats. Bei einem Einsatz von 2.000\,€ muss sich der Index folglich um
etwa 0{,}3\,\% in Regimerichtung bewegen, bevor ein Handelstag den Break-even
erreicht -- eine erreichbare, aber keineswegs sichere Schwelle, deren Konsequenzen
die Evaluation quantifiziert. Die Mindestgröße von 1.000\,€ folgt umgekehrt aus der
Gebührenstruktur: Darunter dominiert die Fixkostenlast jede realistische
Tagesbewegung.

Als extern validierte Alternative ist die Regimestrategie \emph{smatrend}
implementiert, die statt eines Kreuzungssignals den Indexstand gegen einen einzelnen
langfristigen Durchschnitt stellt:
\begin{equation}
  R'(t) \;=\;
  \begin{cases}
    \text{LONG}, & \text{falls } P(t) > \mathrm{SMA}_{325}(t),\\
    \text{CASH}, & \text{sonst,}
  \end{cases}
  \label{eq:smatrend}
\end{equation}
und nur bei einem Regimewechsel handelt, ansonsten die Position hält. Beide
Strategien folgen demselben Plugin-Kontrakt (Signalserie, deklarierter Datenbedarf,
Mindesthistorie, Entscheidungsintervall), sodass der geplante A/B-Vergleich auf
getrennten Depots ohne Eingriff in Auditor, Validator oder Broker möglich ist.

\section{Implementierung}
\label{sec:umsetzung-implementierung}

Grundlage jeder Entscheidung ist die Datenschicht, die einem strikten Prinzip folgt:
Gespeichert werden Rohticks, aus denen Bars abgeleitet werden -- nie umgekehrt.
Damit bleibt jede Verdichtung reproduzierbar, und Auffälligkeiten lassen sich bis
zum einzelnen Tick zurückverfolgen. Ein Sanitätsfilter verwirft unplausible Werte
nicht, sondern markiert sie, sodass die Provenienz erhalten bleibt. Lücken im
Datenstrom werden nach ihrer Ursache klassifiziert; ein empirischer Befund des
Betriebs unterstreicht, warum das nötig ist: Die reguläre Stille des Index
außerhalb der US-Kassasitzung hätte an jedem Montagmorgen fälschlich als
Datenfehler gegolten und den Handel strukturell blockiert. Erst die Unterscheidung
zwischen legitimer Sitzungspause und echtem Verbindungsverlust machte das
Qualitätskriterium betriebstauglich -- ein Beispiel dafür, dass sich die
Schwellwerte eines autonomen Systems nicht am Schreibtisch, sondern nur gegen
reale Betriebsdaten kalibrieren lassen.

Der tägliche Zyklus beginnt mit der Erzeugung eines strukturierten
Entscheidungsdokuments (\emph{DecisionRequest}): Die Strategieauswertung liefert
Regime, Indikatorwerte, Signalalter, Volatilitätskennzahlen, Datenqualitätsmetriken
und einen deterministisch berechneten Kostenkontext (Round-Trip-Fixkosten und deren
prozentuale Last je Größenstufe). Ergänzend setzt sie Prüfmarken
(\emph{Audit-Flags}), von denen ein Teil als harte Ausschlusskriterien definiert ist
-- etwa veraltete Daten, jüngste Verbindungsverluste des Datenstroms oder ein
fehlendes Handelsinstrument. Ein vorgeschaltetes Pre-Gate wertet diese harten
Kriterien vor jedem Modellaufruf aus: Steht das Ergebnis deterministisch fest, wird
die Ablehnung ohne Sprachmodellbeteiligung journalisiert; die Modelle werden nur
befragt, wenn ihre Bewertung das Ergebnis noch verändern kann.

Im Auditschritt erhält jedes der drei konfigurierten Sprachmodelle denselben
DecisionRequest mit einem Systemauftrag als kritischer Risikoprüfer. Die Antwort ist
über ein erzwungenes JSON-Schema strukturiert und wird bei Temperatur null erzeugt;
jedes Modell muss neben der Zustimmung oder Ablehnung eine eigene, unabhängige
Regimeeinschätzung abgeben. Ein fail-safe ausgelegter Validator prüft jede Antwort:
Schemaverletzungen, Zeitüberschreitungen, Widersprüche zwischen eigener Einschätzung
und Strategiesignal sowie Größenüberschreitungen führen zum Veto des betreffenden
Votums; ein nicht erreichbares Modell zählt konservativ als Ablehnung. Die
Einzelvoten werden nach strikter Mehrheit aggregiert, wobei bei Zustimmung die
konservativste (kleinste) Positionsgröße gilt. Steht eine ablehnende Mehrheit
vorzeitig fest, bricht die Befragung ab (\emph{Early-Stop}); eine bereits erreichte
Zustimmungsmehrheit führt hingegen bewusst nicht zum Abbruch, da weitere Voten die
Größe nur verkleinern können. Sämtliche Voten samt Begründungen und Antwortzeiten
werden im Entscheidungsjournal persistiert.

Die Orderausführung ist mehrstufig gegen die typischen Fehlerbilder verteilter
Systeme abgesichert. Ein Frische-Gate verwirft Entscheide, die nicht vom aktuellen
Handelstag stammen. Vor jeder realen Order wird eine \emph{Intent}-Zeile in das
Orderjournal geschrieben, die als eindeutiger Ausführungsmarker dient: Derselbe
Entscheid kann dadurch -- etwa nach einem manuellen Wiederholungslauf oder einer
doppelten Zeitplanung -- kein zweites Mal ausgeführt werden, und selbst ein
Prozessabsturz zwischen Orderaufgabe und Journalisierung bleibt nachvollziehbar.
Scheitert die Kommunikation mit der Bank nach dem Absenden, gleicht eine
Reconciliation den tatsächlichen Zustand über das Orderbuch der Bank ab; jeder
unklare Ausgang endet in einem alarmierten Handelsstopp statt in einem automatischen
Wiederholungsversuch. Einzelorder- und Portfolio-Obergrenzen wirken als
Circuit-Breaker, und die Scharfschaltung auf reale Orders ist ausschließlich
kontospezifisch möglich. Die Bankanbindung selbst nutzt den OAuth-Flow der
comdirect-Schnittstelle mit einer täglich per photoTAN freigegebenen Session-TAN;
der Mensch bleibt über einen Signal-Messenger-Kanal eingebunden, der Tagesreports
liefert und einen jederzeit wirksamen Not-Aus (Pausieren, Fortsetzen, Statusabfrage,
Glattstellen) bereitstellt.

Die Banksitzung selbst folgt einem täglichen Ritual, das die regulatorische
Zwei-Faktor-Anforderung in den autonomen Ablauf integriert: Der morgendliche
Anmeldevorgang fordert genau eine photoTAN-Freigabe an, die der Betreiber
fernbedient über den Messaging-Kanal bestätigt; anschließend hält ein
Erneuerungsdienst die Sitzung über den Handelstag, sodass Orders ohne weitere
Einzelfreigaben, aber stets innerhalb der freigegebenen Sitzung erfolgen. Die
gesamte Broker-Schicht ist zudem kontofähig ausgelegt: Zugangsdaten, Sitzungen,
Journale und der Ausführungsmodus sind je Konto getrennt, sodass der geplante
Parallelbetrieb zweier Depots mit unterschiedlichen Strategien keine strukturelle
Änderung erfordert. Zugangsdaten liegen verschlüsselt in einem dedizierten
Geheimnisspeicher und erscheinen zu keinem Zeitpunkt im Kontext der Sprachmodelle.

Das öffentliche Dashboard schließlich ist als eigenständiges, bewusst kleines
Teilsystem implementiert. Ein Exportprozess redigiert das Entscheidungsjournal auf
die veröffentlichungsfähigen Inhalte, lädt sie nur bei inhaltlicher Änderung hoch
und trennt so den öffentlichen Lesezugriff vollständig vom Handelspfad. Die
Oberfläche stellt den Indexverlauf mit den gleitenden Durchschnitten und den
Regimewechsel-Signalen dar, zeigt je Entscheidung das aggregierte Votum mit
aufklappbaren Einzelbegründungen der Modelle und ist durchgängig responsiv
umgesetzt -- auf Mobilgeräten werden Ordertabellen in vertikale
Beschriftungs-Wert-Karten umgebrochen, sodass die Darstellung ohne horizontales
Scrollen barrierearm lesbar bleibt.

Die Betriebsführung rundet die Implementierung ab. Ein Zeitplandienst startet den
Entscheidungszyklus an Handelstagen kurz vor der US-Markteröffnung; ein
Börsenkalender überspringt Feiertage, und eine Prozesssperre je Konto verhindert,
dass sich zwei Läufe überlappen. Der Not-Aus des Messaging-Kanals ist als
Ein-Buchstaben-Protokoll ausgeführt: Pausieren stoppt jede neue Order vor dem
nächsten Zyklus, Fortsetzen hebt die Sperre wieder auf, die Statusabfrage meldet
den Betriebszustand samt wirksamem Ausführungsmodus, und das Glattstellen schließt
offene Positionen aktiv -- wobei die Pause bewusst erhalten bleibt, bis der Mensch
sie aufhebt. Jeder Zyklus schließt mit einem Bericht an die private
Messaging-Gruppe, der Entscheidung, Voten, Auffälligkeiten und den effektiven
Modus je Konto ausweist; die Scharfschaltung auf reale Orders ist damit täglich
sichtbar dokumentiert.

Die Qualitätssicherung stützt sich auf zwölf Offline-Testsuiten, die von der
Datensammlung über den Voting-Mechanismus bis zu den Schutzpfaden der
Orderausführung jede sicherheitsrelevante Verzweigung mit gemockten Schnittstellen
durchspielen; die Volltexte der zentralen Module sind im Anhang dokumentiert.

\section{Evaluation}
\label{sec:umsetzung-evaluation}

Die Bewertung des Agenten erfolgt auf drei Ebenen: einem historischen Backtest der
Strategie, einer Analyse der realen Transaktionskosten sowie einer laufenden
Kalibrierung des Sprachmodell-Votings im Betrieb.

\subsection{Untersuchungsdesign und Datenbasis}
\label{sec:umsetzung-evaluation-design}

Datenbasis des Backtests sind zehn Jahre echter Tagesschlusskurse des S\&P~500, die
über die Datensammlung des Agenten bezogen und gegen die FRED-Zweitquelle
gegengeprüft wurden. Der Strategie werden drei Vergleichsmaßstäbe gegenübergestellt:
der tägliche Round-Trip der Regel selbst (B0), das passive Halten des Index (B1)
sowie eine Variante, die nur bei einem Regimewechsel handelt (B2). Das Kostenmodell
verwendet keine Annahmen, sondern die real gemessenen Ex-Ante-Kostenausweise der
Bank -- im Direkthandel 7,80\,€ je Round-Trip -- und berücksichtigt die
Kapitalertragsteuer unter Teilfreistellung. Ergänzend misst ein wöchentlicher
Auswertungslauf das Abstimmungsverhalten der Sprachmodelle im laufenden Betrieb: je
Modell die Zustimmungsquote, die Schemafehlerrate, die Antwortlatenz sowie die
Veto-Trefferquote gegen den Folgetags-Return, also den Anteil der Ablehnungen, die
tatsächlich einen Verlusttag in Signalrichtung vermieden hätten, einschließlich des
kontrafaktischen Gewinn- und Verlustbeitrags der Vetos.

Zur Einordnung der Handelsfrequenz dient zusätzlich ein einfaches
Rentabilitätsmodell. Bezeichnet $C(S)$ die Round-Trip-Kosten bei Kapitaleinsatz $S$
und wird die mittlere absolute Intraday-Bewegung des Index mit rund $0{,}6\,\%$
angesetzt, so ergibt sich die für einen positiven Erwartungswert erforderliche
Trefferquote als
\begin{equation}
  p^{*}(S) \;=\; 0{,}5 \;+\; \frac{C(S)}{4 \cdot S \cdot 0{,}006}.
  \label{eq:breakeven}
\end{equation}
Für 1.000\,€ Einsatz liegt diese Schwelle bei etwa 91\,\%, für 2.000\,€ bei etwa
75\,\%; selbst im Grenzfall sehr großer Einsätze setzt der Spread eine
kapitalunabhängige Untergrenze von rund 58\,\% -- Werte, denen die Literatur für
Intraday-Signale einfacher Durchschnittsstrategien lediglich Trefferquoten nahe der
Zufallsschwelle gegenüberstellt. Das Modell begründet damit vorab, weshalb die
tägliche Handelsfrequenz als Demonstrations- und nicht als Renditevehikel zu
verstehen ist.

\subsection{Ergebnisbetrachtung}
\label{sec:umsetzung-evaluation-ergebnisse}

Das zentrale Backtest-Ergebnis fällt eindeutig aus: Der tägliche Round-Trip (B0)
verliert über zehn Jahre 29.598,79\,€, wovon 28.937,94\,€ auf Transaktionskosten
entfallen, bei einer Trefferquote von 31,7\,\%. Die nur bei Regimewechseln handelnde
Variante (B2) erzielt mit neun Positionswechseln in zehn Jahren rund $+100\,\%$ bei
lediglich 106\,€ Kosten; das passive Halten des Index (B1) rund $+250\,\%$. Die
tägliche Handelsfrequenz wird demnach von den Fixkosten aufgezehrt -- die
Rentabilitätsanalyse zeigt, dass bei realistischen Trefferquoten kein
Kapitaleinsatz die tägliche Round-Trip-Strategie in den positiven Erwartungswert
hebt. Diese Erwartung wurde bewusst ehrlich formuliert: Das Erkenntnisziel der
Arbeit ist nicht eine positive Handelsrendite, sondern die autonome, abgesicherte
Ausführung selbst sowie der messbare Beitrag des Sprachmodell-Gremiums; die
Differenz zwischen B0 und B2 beziffert insofern den Preis der täglich demonstrierten
Autonomie. Die Voting-Kalibrierung aus dem Live-Betrieb -- Trockenlauf seit dem
21.07.2026, Echtgeldbetrieb ab dem 27.07.2026 -- ergänzt diese Betrachtung
fortlaufend um die Frage, ob die Vetos der Modelle systematisch Verlusttage treffen
oder lediglich Rauschen darstellen.

Die Kostenmessung gegen das Produktivsystem der Bank lieferte darüber hinaus
eigenständige Befunde. Die Kaufprovision erwies sich über alle getesteten
Ordergrößen als fixer Betrag von 3{,}90\,€, womit sich die Round-Trip-Fixkosten von
7{,}80\,€ bestätigen; die Wahl des Direkthandels beim Emittenten statt eines
Börsenplatzes senkt die Gesamtkosten je Round-Trip um weitere 5\,€, da die
Handelsplatzentgelte entfallen. Zugleich zeigte sich ein versteckter Kostenfaktor:
Der Geld-Brief-Spread des Direkthandels wird im Ex-Ante-Kostenausweis nicht
ausgewiesen und muss im Betrieb empirisch gemessen werden. Auch der Betrieb der
Sprachmodelle erbrachte belastbare Beobachtungen: Die Schema-Adhärenz kleinerer
gehosteter Modelle musste schrittweise erzwungen werden -- von anfänglich
verworfenen Antworten bis zur stabilen, schema-konformen Ausgabe --, für
Modelle mit aktivierter Reasoning-Ausgabe war deren Abschaltung erforderlich, weil
der Denkmodus das erzwungene Antwortformat überstimmte, und ein zunächst
vorgesehenes viertes Modell wurde nach Verifikationsläufen wegen langer
Antwortzeiten und inhaltlichen Rauschens bewusst ausgeschlossen. Die
Antwortlatenzen fielen nach der Warmlaufphase von knapp zehn Sekunden auf unter
eine Sekunde je Votum.

\subsection{Limitationen}
\label{sec:umsetzung-evaluation-limitationen}

Die Aussagekraft der Ergebnisse unterliegt mehreren Einschränkungen. Die
Bankschnittstelle bietet keine Sandbox, sodass die Ordervalidierung gegen das
Produktivsystem erfolgt und die Ausführungsseite bis zum Echtgeldbetrieb nur bis zur
Kostenvorschau demonstriert werden konnte. Die Autonomie ist regulatorisch durch die
täglich erforderliche photoTAN-Freigabe der Banksitzung begrenzt. Im verteilten
Aufbau beträgt die Datenfrische am Entscheidungspunkt 60 bis 90 Sekunden, weshalb
ausschließlich Limit-Orders verwendet werden. Der Backtest approximiert das
Faktor-Zertifikat über das zweifache Index-Tagessegment und modelliert weder
Währungseffekte noch den variablen Geld-Brief-Spread exakt; die inoffizielle
Kursdatenquelle bietet keine Verfügbarkeitsgarantie. Trotz Temperatur null und
erzwungenem Antwortschema bleibt ein Restmaß an Nichtdeterminismus der
Sprachmodelle bestehen, und die Stichprobe der Voting-Kalibrierung ist zu Beginn
des Betriebs klein, weshalb die Auswertung bewusst deskriptiv bleibt und keine
automatische Modellabschaltung auslöst.

\section{Kritische Würdigung}
\label{sec:umsetzung-wuerdigung}

Autonomer Handel mit gehebelten Derivaten vereint mehrere Risikoklassen: das
Marktrisiko des gehebelten Instruments, das Betriebsrisiko eines unbeaufsichtigt
laufenden Systems und das spezifische Risiko probabilistischer Modellkomponenten.
Die Architektur begegnet dem mit einer gestaffelten Sicherungskette, deren
tragendes Prinzip die begrenzte Diskretion ist: Das Sprachmodell-Gremium kann
Handelsvorschläge verhindern, verkleinern oder zeitlich einordnen, niemals aber
erzeugen oder in ihrer Richtung umkehren, und deterministische Ausschlusskriterien
werden grundsätzlich vor jedem Modellaufruf entschieden. Auf der Ausführungsseite
verhindern das zweiphasige Intent-Journal die Doppelausführung, die
Orderbuch-Reconciliation unbemerkte Zustandslücken nach Kommunikationsfehlern und
die Einzel- sowie Portfolio-Obergrenzen ein unkontrolliertes Anwachsen der
Positionen; jeder unklare Zustand mündet in einem alarmierten Handelsstopp. Der
Mensch bleibt dreifach eingebunden: durch die tägliche photoTAN-Freigabe als
regulatorisch erzwungenen Kontrollpunkt, durch den in Echtzeit wirksamen Not-Aus
über den Messaging-Kanal und durch einen Totmannschalter, der den Ausfall des
Agenten selbst meldet. Das öffentliche Dashboard ergänzt die Kontrolle um
Transparenz, ohne die handelnden Systeme zu exponieren.

Gleichwohl bleiben Restrisiken bestehen. Der Ausstiegspfad bei einem Regimewechsel
ist bis zum geplanten Soll-Positions-Abgleich bewusst auf einen alarmierten Stopp
mit manueller Behandlung reduziert -- eine konservative Zwischenlösung, die
menschliches Eingreifen erfordert, statt fehlerhaft zu automatisieren. Die
Wirksamkeit des Modell-Gremiums ist zu Betriebsbeginn empirisch noch unbelegt und
wird erst durch die laufende Kalibrierung überprüfbar; korrelierte Fehlurteile der
Modelle bleiben trotz Anbieterdiversität möglich, da alle Modelle denselben
Entscheidungskontext erhalten. Wirtschaftlich ist festzuhalten, dass die täglich
handelnde Demonstrationsstrategie nach Kosten strukturell defizitär ist und der
Agent insofern ein Erkenntnis-, kein Renditeinstrument darstellt. Von einer
Anlageberatung grenzt sich die Arbeit ausdrücklich ab: Der Agent handelt
ausschließlich für das eigene Experimentalkonto des Betreibers, und weder die
Strategie noch die Systemarchitektur stellen eine Empfehlung an Dritte dar.
Schließlich zeigte die Umsetzung einen strukturellen Befund des Marktumfelds: Die
kostengünstigsten Ausführungswege für Privatanleger sind systematisch nicht über
offene Schnittstellen zugänglich, sodass autonome Agenten einen faktischen
\enquote{API-Aufpreis} tragen -- ein Umstand, der die angestrebte, weitgehend
kostenfreie lokale Betreibbarkeit umso relevanter macht, da sie zumindest die
Infrastruktur- und Inferenzkosten eliminiert.

Über den konkreten Anwendungsfall hinaus beansprucht die Architektur
Übertragbarkeit. Strategie, Produktzuordnung, Modellgremium und Bankadapter sind
über definierte Kontrakte entkoppelt, sodass sich einzelne Bausteine austauschen
lassen, ohne die Sicherungskette zu berühren: Eine andere Regelstrategie ist ein
weiteres Plugin, ein anderes Sprachmodell eine Konfigurationszeile, ein anderer
Broker ein neuer Adapter hinter derselben Orderabstraktion. Für die geplante
Migration auf einen Mac~mini verbleiben als offene Arbeiten im Kern die Anpassung
der Zeitbudgets an langsamere lokale Inferenz, der Ersatz einzelner
Linux-spezifischer Betriebswerkzeuge und die Umstellung der Zeitplanung -- der
Python-Kern selbst ist frei von Cloud-Abhängigkeiten gehalten. Damit bleibt das
Versprechen der Zielsetzung, dieselbe Architektur vom Cloud-Prototyp in einen
privat betreibbaren, laufend kostenfreien Betrieb zu überführen, strukturell
eingelöst, auch wenn der empirische Nachweis dieser letzten Stufe aussteht.

%  =============================================================================
%  Ende
%  =============================================================================
